% from main.tex
%%%%%%%%%%  scan idx 285  |  folio 272  %%%%%%%%%%
\emph{Preuve.} D'après (4.8), \(Rf_*(E)\) est un complexe parfait. L'assertion (i) (resp. (ii)) résulte donc de (5.5 (i)) (resp. (5.6 (i)) ). Enfin, si \(M\) est un \(\mathcal O_Y\)-Module quasi-cohérent, la flèche canonique
\[
 Rf_*(E)\otimes_Y^{\mathrm L}M\longrightarrow Rf_*(E\otimes_Y^{\mathrm L}M)
\]
est un isomorphisme d'après (3.7), donc l'assertion (iii) résulte de (5.6 (ii)).

\textbf{Remarques 5.8.}
\begin{enumerate}[label=\alph*)]
\item Dans la situation de (5.7), notons \(i_y:X_y\to X\) l'injection de la fibre de \(X\) au-dessus d'un point \(y\) de \(Y\). On peut montrer (cf. (EGA III 6 et (IV 3.1.0) que pour \(f\) plat (resp. \(E\) borné à composantes quasi-cohérentes et \(Y\)-plates) on a la formule de Künneth :
\[
 H^i(Rf_*(E)\otimes_Y^{\mathrm L}k(y))\xrightarrow{\sim}H^i(X_y,E_y),
\]
où \(E_y=Li_y^*(E)\) (resp. \(i_y^*(E)\)).
\item On
